[Y=α,W=∃α]: sα(α) = −α ⟶  Hol⁡(γ) = −1 , W = 1 ≜ ΔθL = L ⟶ ω = L⋅f ≡ α = ω

1. Field Realization: Continuous embedding forces a closed trajectory of winding number \(W=1\)

Start from the rigid algebraic core of the YA!WAE! chain. Complete self-reflection is expressed by the order-2 Weyl reflection
\[
s_\alpha(\alpha)=-\alpha.
\]
This operator is an involution: applying it twice recovers the identity. When the discrete group \(\{1,s_\alpha\}\simeq\mathbb{Z}/2\mathbb{Z}\) is realized geometrically as the holonomy of a flat connection (or as monodromy) on a circle, the image of the non-identity element must be a point of order exactly two. Under the standard embedding \(\mathbb{Z}/2\mathbb{Z}\hookrightarrow S^1\) that point is \(-1\), so
\[
\operatorname{Hol}(\gamma)=e^{i\phi_N}=-1
\]
(with \(\phi_N\) any odd multiple of \(\pi\)).

To place this algebraic fact inside a continuous geometric field one must embed the circle into the Euclidean plane (or, equivalently, work on the punctured plane \(M=\mathbb{R}^2\setminus\{(0,0)\}\)). The only continuous trajectories that can carry an order-2 holonomy about the puncture are closed curves that enclose the origin. By definition of the fundamental group
\[
\pi_1(M)\simeq\mathbb{Z},
\]
every closed path is classified by an integer winding number \(W\in\mathbb{Z}\). The generator of this group is any simple closed curve that encircles the origin once in the positive direction; its winding number is \(W=+1\).

Thus the continuous realization of the discrete reflection forces three simultaneous conditions:

the path \(\gamma\) must be closed (otherwise the holonomy is not defined as a group element of the monodromy representation);
the path must enclose the puncture (otherwise the holonomy is trivial);
the oriented enclosure must be the fundamental generator, hence \(W=1\).

No further metric structure is required for this step. The forcing is purely topological: the algebraic requirement that the holonomy be non-trivial of order two selects the non-zero element of \(\pi_1(M)\), and the simplest positive representative of that class has winding number one. Any other integer winding number \(W=k\) would correspond to a \(k\)-fold cover of the same reflection; the minimal, irreducible realization of the order-2 element itself is the generator \(W=1\).

(The pure Cartesian square constructed without trigonometric functions or any occurrence of \(\pi\) demonstrates the same fact: successive oriented turns accumulate to the integer \(+1\), confirming that the winding number is independent of any particular angular chart.)

2. Final Mathematical Closure: \(\alpha=\omega\) and the coordinate-language dependence of the numerical coefficient

Once a continuous trajectory of winding number \(W=1\) has been forced, the angular rate at which that trajectory is traversed becomes the dynamical quantity \(\omega\). By definition the angular frequency is the amount of angular coordinate space swept out per unit time. If the system completes \(f\) full cycles per unit time, then
\[
\omega=\bigl(\text{angular measure of one cycle}\bigr)\times f.
\]

The angular measure of one cycle is not an absolute number; it is the length assigned to the generator of \(\pi_1(S^1)\) by the chosen coordinate chart on the circle:

In the standard radian chart the generator is normalized to length \(2\pi\), so
  \[
  \omega=2\pi f.
  \]
In the pure-turn (cycle) chart the generator is normalized to length \(1\), so
  \[
  \omega=1\cdot f=f.
  \]
In the gradian chart the generator is normalized to length \(400\), so \(\omega=400f\), and so on.

Because the original algebraic object being reflected is the simple root \(\alpha\) itself, and because the continuous trajectory is the geometric realization of that same reflection, the identification
\[
\alpha=\omega
\]
holds at the level of the structural correspondence: the generator of the discrete reflection is identified with the generator of the continuous angular motion. The numerical coefficient that multiplies the linear frequency \(f\) is merely the chart-dependent length of that generator. Changing the chart changes the number; it does not change the underlying bijection between the algebraic reflection and the continuous cycle.

Consequently the governing identity
\[
[Y=\alpha,\,W=\exists!\alpha]\;\longrightarrow\;\operatorname{Hol}(\gamma)=-1\;\longrightarrow\;W=1\;\longrightarrow\;\alpha=\omega
\]
remains intact under every admissible angular coordinate system. Only the concrete real number that appears when \(\omega\) is written in a particular language is chart-dependent. The structural core—order-2 reflection realized as a once-wound continuous trajectory whose generator is identified with the original root—is independent of that language and is therefore unassailable by a change of units.

In short:

Topology forces the closed trajectory of winding number \(W=1\).
The angular frequency \(\omega\) is the angular measure of that trajectory per unit time.
The numerical prefactor is the measure of one full cycle in the chosen chart.
The identification \(\alpha=\omega\) is the chart-independent structural equality between the algebraic generator and the continuous generator.

This completes the deterministic chain from discrete self-reflection to continuous dynamical frequency without inserting any privileged numerical constant by hand.

Thorough detailing of the gradian (and general angular-unit) case

Angular charts as different normalizations of the same generator

The fundamental group of the circle (or of the punctured plane) is
\[
\pi_1(S^1)\simeq\mathbb{Z}.
\]
Its generator is any simple closed curve that winds once about the origin in the positive direction. This generator is a purely topological object; it carries no intrinsic numerical length. A numerical length appears only when one chooses a continuous coordinate chart that assigns a positive real number to that generator. Different charts simply choose different numbers for the same topological object.

Three common charts illustrate the pattern:

Radian chart**  
  The generator is assigned length \(2\pi\).  
  One full cycle therefore contributes \(2\pi\) units of angular coordinate.  
  If the system completes \(f\) cycles per unit time,
  \[
  \omega_{\text{rad}}=2\pi\,f.
  \]

Turn (cycle) chart**  
  The generator is assigned length \(1\).  
  One full cycle contributes \(1\) unit.  
  \[
  \omega_{\text{turn}}=1\cdot f=f.
  \]

Gradian chart**  
  The generator is assigned length \(400\).  
  (Historically, the gradian was defined so that a right angle is \(100\) gradians and a full circle is therefore \(400\) gradians.)  
  One full cycle contributes \(400\) units of angular coordinate.  
  Consequently the angular frequency measured in gradians per unit time is
  \[
  \omega_{\text{grad}}=400\,f.
  \]

Any other positive real number \(L\) can serve as the length of the generator. The corresponding frequency is simply
\[
\omega_L=L\cdot f.
\]
The family of all such charts is parametrized by the single free positive real \(L=\text{“angular measure of one topological cycle”}\).

Why the structural identification \(\alpha=\omega\) survives every chart

Recall the logical chain:

Complete self-reflection forces the algebraic identity \(s_\alpha(\alpha)=-\alpha\), an element of order two.
Continuous realization of that order-two element as holonomy on the circle selects the non-trivial generator of \(\pi_1(S^1)\), i.e., a trajectory of winding number \(W=1\).
The angular frequency \(\omega\) is defined as the rate at which this generator is traversed.
Because the object being reflected is the simple root \(\alpha\) itself, the generator of the discrete reflection is identified with the generator of the continuous motion: \(\alpha=\omega\).

None of these four steps refers to a particular numerical length \(L\). The equality \(\alpha=\omega\) is therefore an equality of generators, not an equality of real numbers. When a concrete chart is chosen, both sides of the equality are expressed in the same units and the numerical coefficient \(L\) appears on the right-hand side:
\[
\alpha\;\longleftrightarrow\;\omega_L=L\cdot f.
\]
Changing the chart merely multiplies both the angular measure and the reported frequency by the same constant; the underlying bijection between the algebraic reflection and the continuous cycle remains unchanged.

Explicit translation between charts

Let \(\theta_L\) be the continuous angular coordinate in the chart that assigns length \(L\) to the generator. Then the change in that coordinate after one positive winding is exactly \(L\). The conversion rules are therefore linear:

\[
\begin{align*}
\theta_{\text{rad}}&=2\pi\cdot\theta_{\text{turn}},\\
\theta_{\text{grad}}&=400\cdot\theta_{\text{turn}},\\
\theta_{\text{grad}}&=\frac{400}{2\pi}\cdot\theta_{\text{rad}},
\end{align*}
\]
and likewise for the corresponding frequencies. Because every chart is related to every other by a positive constant factor, the dynamical content of the YA!WAE! identity is invariant under the entire family of charts.

Summary of the gradian case inside the larger thesis

Topology forces a closed trajectory of winding number \(W=1\).
The gradian chart simply declares that this trajectory has angular measure \(400\).
Therefore the angular frequency expressed in gradians is \(\omega=400f\).
The identification \(\alpha=\omega\) continues to hold because both sides refer to the same generator; only the numerical label attached to that generator has changed.
The same reasoning applies to every other admissible angular unit. The structural core of the equation is consequently independent of the particular language (radians, turns, gradians, \ldots) in which the angular measure is written.

The numeric coefficient is always the chart-dependent length of one topological cycle; the forcing of the cycle itself, and the identification of that cycle with the original simple root, remain chart-independent.

The simple root \(\alpha\) is not a spectator.  
It is the very thing that turns upon itself.  
When the mirror of reflection is held up,  
\(\alpha\) meets its own negative and recognizes  
the only possible answer: an order-two flip.

That same \(\alpha\), once released from the discrete lattice  
and allowed to walk in continuous space,  
must travel a closed path that returns it to itself  
after a single, irreducible circuit.  
The path is not an external stage;  
it is the geometric body of the reflection itself.

Therefore the rate at which this path is traversed  
—the angular frequency \(\omega\)—  
is not a separate actor.  
It is \(\alpha\) in motion.  
The generator that flipped in algebra  
is the generator that turns in geometry.

So the equation \(\alpha=\omega\) is not a forced marriage of two strangers.  
It is the quiet recognition that the thing being reflected  
and the thing that is turning  
were always true reflections mathematically, topologically, geometrically and beyond.

This is the chart-independent skeleton:

The initial pair already fixes a unique simple root and its order-two reflection.  
Continuous realization forces the holonomy \(-1\) and winding number \(W=1\).  
Continuous argument along that path equals the chart-dependent length \(L\) of the generator.  
Angular frequency is that length traversed at linear rate \(f\).  
Because the moving object is the original root itself, the algebraic generator and the dynamical generator coincide: \(\alpha\equiv\omega\).

Only the numerical value of \(L\) changes with the angular language; the structural forcing remains absolute.
1. The initial datum and what it already contains

The pair
\[
[Y=\alpha,\quad W=\exists!\alpha]
\]
is not an arbitrary starting point. It encodes a single, rigid geometric fact drawn from the theory of root systems and Weyl groups:

A fundamental chamber has been fixed.
The walls of that chamber are in simple transitive correspondence with the simple roots.
Therefore, for any chosen wall there exists a unique simple root \(\alpha\).
The reflection through the hyperplane orthogonal to that root satisfies the universal identity
  \[
  s_\alpha(\alpha)=-\alpha.
  \]

This identity is an involution of order exactly two. No further choice remains at the algebraic level: the object, the uniqueness of the object, and the order-two character of its self-action are already fixed by the notation \([Y=\alpha,\,W=\exists!\alpha]\).

2. Forced topological realization

Once the algebraic involution is required to act continuously, it must be realized as the holonomy (or monodromy) of a closed path \(\gamma\). The only continuous image of an order-two element inside \(S^1\) is a point of order two, conventionally written
\[
\operatorname{Hol}(\gamma)=e^{i\phi_N}=-1.
\]
Because the base space is the punctured plane (or any space homotopy-equivalent to a circle), the fundamental group is \(\mathbb{Z}\). The non-trivial order-two holonomy selects a non-zero winding class; the generator of that class is the simple closed curve of winding number
\[
W=+1.
\]
Thus the moment the path is required to carry the holonomy of the original reflection, both the closed character of the trajectory and its winding number are forced. No other integer satisfies the minimal, irreducible realization of the order-two element.

3. Forced continuous argument accumulation

Along any such path the continuous determination of the argument function (the unique continuous lift of \(\operatorname{atan2}\) after analytic continuation or numerical unwrapping) changes by exactly one full generator of the circle. In an arbitrary angular chart that assigns length \(L\) to the generator, the measured accumulation is
\[
\Delta\theta_L=L.
\]
In the radian chart \(L=2\pi\); in the turn chart \(L=1\); in the gradian chart \(L=400\); and so on. The accumulation itself is forced by the topology; only the numerical label attached to it is chart-dependent.

4. Forced metric identification of the period and the frequency

The intrinsic period \(\tau\) of the cycle is identified with the measured accumulation:
\[
\tau\equiv\Delta\theta_L=L.
\]
The angular frequency is the rate at which this period is traversed:
\[
\omega=L\cdot f.
\]
Because the original algebraic object being reflected is \(\alpha\) itself, and because the continuous trajectory is the geometric body of that same reflection, the generator of the discrete involution is identical with the generator of the continuous motion. This yields the structural equality
\[
\alpha=\omega
\]
(understood as an equality of generators, independent of the particular numerical value of \(L\)).

Once the initial path realizes the unique simple root and its order-two reflection as continuous holonomy, every subsequent link is compelled by topology, by the definition of continuous argument, and by the identification of the object with its own geometric motion. No additional free parameter is introduced; the only remaining freedom is the conventional choice of angular chart, which merely rescales the numerical coefficient \(L\) while leaving the structural skeleton untouched.

6. Summary statement of the thesis

The notation \([Y=\alpha,\,W=\exists!\alpha]\) already contains a unique simple root and an order-two reflection. Requiring that reflection to act continuously forces a closed trajectory of winding number one. The continuous argument along that trajectory forces a period equal to the chart-dependent length of the generator. The angular frequency is that period traversed at linear rate \(f\). Because the moving object is the original root itself, the algebraic generator and the dynamical generator coincide: \(\alpha=\omega\).  

Thus the entire architecture—topological, geometric, metric, and dynamical—is forced by the single initial decision to let the unique simple root follow its own reflective path into continuous space.

The numerical coefficient \(L\) is the sole chart-dependent quantity ( \(L=2\pi\) in radians, \(L=1\) in turns, \(L=400\) in gradians, \ldots). Every other arrow is forced once the initial pair \([Y=\alpha,\,W=\exists!\alpha]\) is required to act continuously.
